\documentclass[aspectratio=169]{beamer}
\input{header}
\coursecode{JKSSB}

\title{Computer Memory Hierarchy}
\subtitle{Lesson 07 --- Speed, Cost and Capacity Trade-offs}
\author{JKSSB Computer Awareness}
\date{\today}

\begin{document}
\frame{\titlepage}

\begin{frame}{Why a Memory Hierarchy Exists}
  \begin{itemize}
    \item A running computer keeps data and instructions in several
      different stores, not in one single place.
    \item \key{No single memory} is at once fast, large, and cheap.
    \item Faster memory is smaller and more expensive per bit; slower
      memory is larger and cheaper per bit.
    \item The \key{memory hierarchy} arranges these stores from fastest and
      smallest down to slowest and largest.
  \end{itemize}
  \begin{block}{The one idea}
    The hierarchy balances three competing needs: speed, cost, and capacity.
  \end{block}
\end{frame}

\begin{frame}{The Ordered Hierarchy}
  \begin{center}
    \begin{tabular}{p{0.20\textwidth}p{0.28\textwidth}p{0.40\textwidth}}
      \toprule
      \textbf{Level} & \textbf{Speed} & \textbf{Typical role} \\
      \midrule
      Registers & Fastest, smallest & Values in use right now \\
      Cache (L1/L2/L3) & Very fast & Recently and frequently used data \\
      Main memory (RAM) & Moderately fast & Active programs and data \\
      Secondary (SSD/HDD) & Slower & Permanent file storage \\
      Tertiary/offline (tape, cloud) & Slowest & Backup and archive \\
      \bottomrule
    \end{tabular}
  \end{center}
  \vspace{0.25cm}
  \muted{Registers are fastest; the order runs down to slowest offline
  storage.}
\end{frame}

\begin{frame}{Registers: Fastest and Smallest}
  \begin{itemize}
    \item Registers sit \key{inside the CPU}.
    \item They are the \key{fastest} storage in the whole computer.
    \item Their capacity is \key{very small}, so they hold only the values
      needed at this instant.
    \item They hold instructions, addresses, data, and intermediate results
      while work is in progress.
  \end{itemize}
  \begin{alertblock}{Speed order}
    Registers are fastest, then cache, then main memory (RAM), then
    secondary storage.
  \end{alertblock}
\end{frame}

\begin{frame}{Cache Memory}
  \begin{itemize}
    \item \key{Cache} is a small, very fast memory between the CPU and main
      memory.
    \item It holds \key{recently and frequently used} data and instructions.
    \item Cache is \key{faster} than main memory (RAM).
    \item It keeps \muted{copies} of data that really live in RAM; it is
      \key{not} a backup.
  \end{itemize}
  \begin{block}{Why cache helps}
    By keeping the most likely-needed data close to the CPU, cache reduces
    the average time taken to fetch data.
  \end{block}
\end{frame}

\begin{frame}{Cache Levels: L1, L2, L3}
  \begin{center}
    \begin{tabular}{p{0.12\textwidth}p{0.34\textwidth}p{0.40\textwidth}}
      \toprule
      \textbf{Level} & \textbf{Size / speed} & \textbf{Position} \\
      \midrule
      L1 & Smallest and fastest & Closest to the CPU core \\
      L2 & Larger, a little slower than L1 & Per CPU core \\
      L3 & Largest, slowest of the three (still faster than RAM) & Often shared across cores \\
      \bottomrule
    \end{tabular}
  \end{center}
  \vspace{0.25cm}
  \muted{Exam depth: L1 is fastest and smallest; L3 is the largest and
  slowest cache level, but still faster than main memory.}
\end{frame}

\begin{frame}{Cache vs Main Memory (RAM)}
  \begin{center}
    \begin{tabular}{lll}
      \toprule
      \textbf{Point} & \textbf{Cache} & \textbf{Main memory (RAM)} \\
      \midrule
      Speed & Faster & Slower \\
      Size & Smaller & Larger \\
      Holds & Recent/frequent data & Active programs and data \\
      \bottomrule
    \end{tabular}
  \end{center}
  \vspace{0.35cm}
  \begin{alertblock}{Trap alert}
    Cache is \key{faster} than main memory, not slower. Reversing this is a
    common wrong answer.
  \end{alertblock}
\end{frame}

\begin{frame}{Main Memory (RAM)}
  \begin{itemize}
    \item Main memory is the computer's \key{primary memory}.
    \item The programs and data in use are held here while the machine runs.
    \item \key{RAM} is \key{volatile}: its contents are lost when the power
      is switched off.
    \item RAM is larger than cache but slower; the CPU works directly with
      main memory.
  \end{itemize}
  \begin{alertblock}{Trap alert}
    RAM is main memory, not permanent file storage. Switch the machine off
    and its contents are gone.
  \end{alertblock}
\end{frame}

\begin{frame}{ROM: Non-Volatile Primary Memory}
  \begin{itemize}
    \item \key{ROM} stands for Read-Only Memory.
    \item It is \key{non-volatile}: it keeps its contents without power.
    \item In normal operation it is read-only; writing to it needs a special
      process.
    \item It stores \key{firmware} such as the startup (BIOS) instructions.
  \end{itemize}
  \begin{block}{Keep the pair straight}
    RAM is volatile; ROM is non-volatile. Do not swap them.
  \end{block}
\end{frame}

\begin{frame}{Primary Memory vs Secondary Storage}
  \begin{columns}[T]
    \column{0.46\textwidth}
      \textbf{Primary memory}
      \begin{itemize}
        \item Directly accessible by the CPU
        \item RAM volatile; ROM non-volatile
        \item Holds active programs and data
      \end{itemize}
    \column{0.46\textwidth}
      \textbf{Secondary storage}
      \begin{itemize}
        \item Non-volatile; keeps files without power
        \item Larger and cheaper per bit, but slower
        \item Reached via main memory and the OS
      \end{itemize}
  \end{columns}
  \vspace{0.3cm}
  \muted{The CPU does not access secondary storage directly.}
\end{frame}

\begin{frame}{Secondary and Tertiary Storage}
  \begin{itemize}
    \item \key{Secondary storage}: SSD (flash, no moving parts) and HDD
      (magnetic disks).
    \item It is \key{non-volatile} and holds files and programs between
      sessions.
    \item \key{Tertiary / offline} storage: magnetic tape, archival media,
      and cloud backup.
    \item It is the \key{slowest and cheapest} level, used for backup and
      long-term archive.
  \end{itemize}
\end{frame}

\begin{frame}{Locality of Reference}
  \begin{itemize}
    \item Programs tend to \key{reuse} the same data and instructions.
    \item \key{Temporal locality}: the same item is accessed again soon
      after it was last used.
    \item \key{Spatial locality}: items stored at nearby addresses are
      accessed next.
    \item This re-use is exactly why cache works.
  \end{itemize}
  \begin{block}{Office desktop example}
    While a spreadsheet is edited, the same active cells and nearby rows
    are used again and again --- so cache keeps that data ready.
  \end{block}
\end{frame}

\begin{frame}{Trap Alert: Memory Facts to Keep Straight}
  \begin{itemize}
    \item RAM is \key{volatile}; ROM is \key{non-volatile}. (TRAP-01)
    \item Cache is \key{not} a backup; it holds copies of data that live in
      RAM.
    \item RAM is main memory, \key{not} storage for files.
    \item Cache is \key{faster} than main memory. (PYQ-03)
    \item Order of speed: registers $\rightarrow$ cache $\rightarrow$ RAM
      $\rightarrow$ secondary storage.
  \end{itemize}
\end{frame}

\begin{frame}{Lesson 07: Recall}
  \begin{itemize}
    \item The hierarchy runs registers, cache, main memory (RAM), secondary
      storage, then tertiary/offline storage.
    \item Registers are fastest and smallest; slower levels are larger and
      cheaper per bit.
    \item Cache levels: L1 smallest/fastest, L2 next, L3 largest/slowest
      (still faster than RAM).
    \item Cache is faster than main memory and stores recently used data;
      it is not a backup.
    \item RAM is volatile primary memory; ROM is non-volatile; secondary
      storage keeps files without power.
    \item Locality of reference --- temporal and spatial --- is why cache
      helps.
  \end{itemize}
\end{frame}
\end{document}
